\exercisesection{\S~3}

\begin{enumerate}
\def\labelenumi{\arabic{enumi})}
\item
  Let \(\Gamma\) be a subgroup of \(\mathrm{Q}(\mathrm{R})\) of finite index, and \(\mathrm{P} = \Gamma \cap \mathrm{R}\) . The algebra \(\mathfrak{h} + \mathfrak{g}^{\mathrm{P}}\) is reductive, and any reductive subalgebra of \(\mathfrak{g}\) containing \(\mathfrak{h}\) can be obtained in this way (use Chap. VI, §1, Exerc. 6 b)).
\item
  Let X be the set of reductive subalgebras of g, distinct from g, and containing h. Determine the minimal elements of X by means of Chap. VI, §4, Exerc. 4. Show that the centre of such a maximal subalgebra is of dimension 0 or 1, according to whether we are in case a) or case b) of the exercise in question.
\item
  Let \(\mathfrak{b}\) be a Borel subalgebra of \(\mathfrak{g}\) , and let \(l = \mathrm{rk}(\mathfrak{g})\) . Show that the minimum number of generators of the algebra \(\mathfrak{b}\) is \(l\) if \(l \neq 1\) , and 2 if \(l = 1\) .
\item
  Assume that \(k = \mathbf{R}\) or \(\mathbf{C}\) . Put \(G = \operatorname{Int}(\mathfrak{g})\) , and identify the Lie algebra of \(G\) with \(\mathfrak{g}\) . Let \(\mathfrak{b}\) be a Borel subalgebra of \((\mathfrak{g},\mathfrak{h})\) , and \(\mathfrak{n}\) the set of its nilpotent elements. Denote by \(H,B,N\) the integral subgroups of \(G\) with Lie algebras \(\mathfrak{h},\mathfrak{b},\mathfrak{n}\) , respectively. Show that \(H,B,N\) are Lie subgroups of \(G\) , that \(N\) is simply-connected, and that \(B\) is the semi-direct product of \(H\) by \(N\) .
\item
  Let \(\mathfrak{m}\) be a parabolic subalgebra of a semi-simple Lie algebra \(\mathfrak{a}\) .
\end{enumerate}

\begin{enumerate}
\def\labelenumi{\alph{enumi})}
\item
  Let p be a subalgebra of m. Then p is a parabolic subalgebra of a if and only if p contains the radical r of m and p/r is a parabolic subalgebra of the semi-simple algebra m/r.
\item
  If \(m'\) is a parabolic subalgebra of a, every Cartan subalgebra of \(m \cap m'\) is a Cartan subalgebra of a. (Reduce to the split case, and apply Prop. 10 to the Borel subalgebras contained in m and \(m'\) .)
\end{enumerate}

\begin{enumerate}
\def\labelenumi{\arabic{enumi})}
\setcounter{enumi}{5}
\item
  Two Borel subalgebras of a semi-simple Lie algebra \(\mathfrak{a}\) are said to be opposite if their intersection is a Cartan subalgebra. Show that, if \(\mathfrak{b}\) is a Borel subalgebra of \(\mathfrak{a}\) , and \(\mathfrak{h}\) is a Cartan subalgebra of \(\mathfrak{b}\) , there exists a unique Borel subalgebra of \(\mathfrak{a}\) that is opposite to \(\mathfrak{b}\) and contains \(\mathfrak{h}\) . (Reduce to the split case.)
\item
  Let \(\mathfrak{a}\) be a semi-simple Lie algebra, \(\mathfrak{h}\) a Cartan subalgebra of \(\mathfrak{a}\) , and \(\mathfrak{s}\) a semi-simple subalgebra of \(\mathfrak{a}\) containing \(\mathfrak{h}\) . Show that:
\end{enumerate}

\[
(\mathfrak {a}, \mathfrak {h}) \text {is split} \Longleftrightarrow (\mathfrak {s}, \mathfrak {h}) \text {is split.}
\]

Construct an example in which \(\mathfrak{a}\) is splittable but \(\mathfrak{s}\) is not. (Take \(\mathfrak{a} = \mathfrak{sp}(4, k)\) and \(\mathfrak{s} = \mathfrak{sl}(2, k')\) , where \(k'\) is a quadratic extension of \(k\) .)

\begin{enumerate}
\def\labelenumi{\arabic{enumi})}
\setcounter{enumi}{7}
\tightlist
\item
  Choose a basis of R, and hence a set of positive roots \(R_{+}\) . Assume that R is irreducible. Let \(\tilde{\alpha}\) be the highest root, S the set of roots orthogonal to \(\tilde{\alpha}\) , and \(S_{+} = S \cap R_{+}\) .
\end{enumerate}

\begin{enumerate}
\def\labelenumi{\alph{enumi})}
\tightlist
\item
  Let
\end{enumerate}

\[
\mathfrak {g} ^ {\prime} = \mathfrak {g} ^ {\mathrm{S}} \oplus \mathfrak {h} _ {\mathrm{S}}, \quad \mathfrak {m} _ {+} = \sum_ {\alpha \in \mathrm{S} _ {+}} \mathfrak {g} ^ {\alpha}, \quad \mathfrak {m} _ {-} = \sum_ {\alpha \in \mathrm{S} _ {+}} \mathfrak {g} ^ {- \alpha}.
\]

Then \(\mathfrak{g}'\) is semi-simple and \(\mathfrak{g}' = \mathfrak{h}_{\mathrm{S}} \oplus \mathfrak{m}_{+} \oplus \mathfrak{m}_{-}\) .

\begin{enumerate}
\def\labelenumi{\alph{enumi})}
\setcounter{enumi}{1}
\tightlist
\item
  Put:
\end{enumerate}

\[
\mathfrak {n} _ {+} = \sum_ {\alpha \in \mathrm{R} _ {+}} \mathfrak {g} ^ {\alpha}, \quad \mathfrak {p} = \sum_ {\alpha \in \mathrm{R} _ {+} - \mathrm{S} _ {+}} \mathfrak {g} ^ {\alpha}, \quad \mathfrak {p} _ {0} = \sum_ {\alpha \in \mathrm{R} _ {+} - \mathrm{S} _ {+} - \{\tilde {\alpha} \}} \mathfrak {g} ^ {\alpha}.
\]

Then

\[
\mathfrak {n} _ {+} = \mathfrak {m} _ {+} \oplus \mathfrak {p}, \quad [ \mathfrak {m} _ {+}, \mathfrak {p} _ {0} ] \subset \mathfrak {p} _ {0}, \quad [ \mathfrak {n} _ {+}, \mathfrak {g} ^ {\tilde {\alpha}} ] = 0.
\]

In particular, \(\mathfrak{p}\) is an ideal of \(\mathfrak{n}_{+}\) .

\begin{enumerate}
\def\labelenumi{\alph{enumi})}
\setcounter{enumi}{2}
\tightlist
\item
  For all \(\alpha \in \mathrm{R}_{+} - \mathrm{S}_{+} - \{\tilde{\alpha}\}\) , there exists a unique \(\alpha' \in \mathrm{R}_{+} - \mathrm{S}_{+} - \{\tilde{\alpha}\}\) such that \(\alpha + \alpha' = \tilde{\alpha}\) . For all \(\alpha \in \mathrm{R}_{+} - \mathrm{S}_{+} - \{\tilde{\alpha}\}\) , a non-zero element \(X_{\alpha}\) of \(\mathfrak{g}^{\alpha}\) can be chosen so that
\end{enumerate}

\[
\begin{array}{r l} {[ X _ {\alpha}, X _ {\alpha^ {\prime}} ] = \pm X _ {\tilde {\alpha}}} & {\mathrm{if} \alpha \in \mathrm{R} _ {+} - \mathrm{S} _ {+} - \{\tilde {\alpha} \}} \\ {[ X _ {\alpha}, X _ {\beta} ] = 0} & {\mathrm{if} \alpha , \beta \in \mathrm{R} _ {+} - \mathrm{S} _ {+}, \beta \neq \alpha^ {\prime}. ^ {1 1}} \end{array}
\]

\footnotetext[11]{This exercise was communicated to us by A. JOSEPH.}

\begin{enumerate}
\def\labelenumi{\arabic{enumi})}
\setcounter{enumi}{8}
\tightlist
\item
  Construct examples of semi-simple Lie algebras \(\mathfrak{a}\) such that:
\end{enumerate}

\begin{enumerate}
\def\labelenumi{\roman{enumi})}
\item
  a has no Borel subalgebra;
\item
  a has a Borel subalgebra, but is not splittable.
\end{enumerate}

¶10) Let a be a semi-simple Lie algebra, and x an element of a. We say that x is diagonalizable if ad x is diagonalizable (Algebra, Chap. VII), in other words, if there exists a basis of a consisting of eigenvectors of ad x.

\begin{enumerate}
\def\labelenumi{\alph{enumi})}
\item
  Let \(\mathfrak{c}\) be a commutative subalgebra of \(\mathfrak{a}\) consisting of diagonalizable elements, and let L be the set of weights of \(\mathfrak{c}\) in the representation \(\mathrm{ad}_{\mathfrak{a}}\) . The set L is a finite subset of \(\mathfrak{c}^*\) containing 0 (unless \(\mathfrak{a} = 0\) ) and \(\mathfrak{a} = \bigoplus_{\lambda \in L} \mathfrak{a}^{\lambda}(\mathfrak{c})\) . Show that there exists a subset M of L such that \(L - \{0\}\) is the disjoint union of M and \(-M\) , and such that \((M + M) \cap L \subset M\) . If M has these properties, put \(\mathfrak{a}^{M} = \bigoplus_{\lambda \in M} \mathfrak{a}^{\lambda}(\mathfrak{c})\) , and \(\mathfrak{p}^{M} = \mathfrak{a}^{0}(\mathfrak{c}) \oplus \mathfrak{a}^{M}\) . The algebra \(\mathfrak{a}^{0}(\mathfrak{c})\) is the commutant of \(\mathfrak{c}\) in \(\mathfrak{a}\) ; it is reductive in \(\mathfrak{a}\) (Chap. VII, §1, no. 5), and its Cartan subalgebras are the Cartan subalgebras of \(\mathfrak{a}\) (Chap. VII, §2, no. 3). Show that \(\mathfrak{p}^{M}\) is a parabolic subalgebra of a, and that \(a^{M}\) is the set of nilpotent elements of the radical of \(p^{M}\) . (Use the fact that c is contained in a Cartan subalgebra of a and reduce, by extension of scalars, to the case in which this subalgebra is splitting.)
\item
  Retain the hypotheses and notations of a), and assume in addition that a is splittable. Show that c is contained in a splitting Cartan subalgebra of a. (Take a splitting Cartan subalgebra h of a contained in \(p^{M}\) ; then there exists a unique Cartan subalgebra \(h'\) of \(\mathfrak{a}^{0}(\mathfrak{c})\) such that h is contained in \(h' \oplus a^{M}\) ; show that \(h'\) is a splitting Cartan subalgebra of a, and that \(h'\) contains c.)
\end{enumerate}

\begin{enumerate}
\def\labelenumi{\arabic{enumi})}
\setcounter{enumi}{10}
\tightlist
\item
  Let \(\mathfrak{a} = \prod \mathfrak{a}_i\) be a finite product of semi-simple Lie algebras. A subalgebra \(\mathfrak{q}\) of \(\mathfrak{a}\) is a parabolic (resp. Borel) subalgebra if and only if it is of the form \(\mathfrak{q} = \prod \mathfrak{q}_i\) where, for all \(i\) , \(\mathfrak{q}_i\) is a parabolic (resp. Borel) subalgebra of \(\mathfrak{a}_i\) .
\end{enumerate}

¶ 12) Let \(k'\) be a finite extension of \(k\) , \(\mathfrak{a}'\) a semi-simple \(k'\) -Lie algebra, and \(\mathfrak{a}\) the underlying \(k\) -Lie algebra. Show that the parabolic (resp. Borel) subalgebras of \(\mathfrak{a}\) are the same as those of \(\mathfrak{a}'\) . (Extend scalars to an algebraic closure of \(k\) , and use Exerc. 11.)

\begin{enumerate}
\def\labelenumi{\arabic{enumi})}
\setcounter{enumi}{12}
\item
  Let p and q be two parabolic subalgebras of a semi-simple Lie algebra a, and let n be the nilpotent radical of p.~Show that \(\mathfrak{m} = (\mathfrak{p} \cap \mathfrak{q}) + \mathfrak{n}\) is a parabolic subalgebra of a. (Reduce to the case where a is split and q is a Borel subalgebra; choose a Cartan subalgebra h contained in \(p \cap q\) , and determine the subset P of the corresponding root system such that \(m = h + g^{P}\) .)
\item
  Retain the notations of Prop. 9.
\end{enumerate}

\begin{enumerate}
\def\labelenumi{\alph{enumi})}
\item
  Let \(\alpha \in \mathrm{B}\) . Show that \(\mathfrak{n} \cap \operatorname{Kerad} X_{\alpha}\) is the direct sum of the \(\mathfrak{g}^{\beta}\) , where \(\beta\) belongs to the set of elements of \(\mathrm{R}_{+}\) such that \(\alpha + \beta \notin \mathrm{R}_{+}\) .
\item
  Deduce that, if g is simple, the centre of n is equal to \(g^{\tilde{\alpha}}\) , where \(\tilde{\alpha}\) is the highest root of \(R_{+}\) . In the general case, the dimension of the centre of n is equal to the number of simple components of g.
\end{enumerate}
% semantic-fragments: legacy-input-seam
\relax{}
